\documentclass[11pt]{article}
\usepackage[margin=1in]{geometry}
\usepackage{amsmath,amssymb,amsthm}
\usepackage{booktabs}
\usepackage{hyperref}
\usepackage{enumitem}
\hypersetup{colorlinks=true,linkcolor=blue,urlcolor=blue}

\newcommand{\QQ}{\mathbb{Q}}
\newcommand{\PP}{\mathbb{P}}
\newcommand{\E}{\mathbb{E}}
\newcommand{\R}{\mathbb{R}}
\newcommand{\dd}{\mathrm{d}}
\newcommand{\sg}{\operatorname{sg}}
\DeclareMathOperator*{\lse}{logsumexp}
\DeclareMathOperator{\Var}{Var}
\DeclareMathOperator{\std}{std}

\title{\textbf{IRM2: A Neural Short-Rate Model for Joint Calibration of\\
Treasury Yield Curves and Futures}\\[2pt]
\large Mathematical specification, diagnostics, and open problems}
\author{Project report}
\date{June 2026}

\begin{document}
\maketitle

\begin{abstract}
IRM2 calibrates a latent-state neural stochastic differential equation (SDE)
to the U.S.\ Treasury yield curve and Treasury futures. A recurrent encoder
maps a window of observed curves to a latent initial condition; a Neural SDE
evolves it under the risk-neutral measure $\QQ$; a decoder produces the short
rate; and a Monte-Carlo pricer values zero-coupon bonds and the
cheapest-to-deliver futures option. This document gives the precise
formulation, the sequence of identifiability problems uncovered during
development (volatility, the term premium, and the decoupling of the futures
head), the terms added to resolve them, and the remaining open questions. It
is written for a reader comfortable with stochastic calculus and is
deliberately concise on engineering detail.
\end{abstract}

\section{State-space construction}
Let $\mathbf c_t\in\R^M$ denote the observed yield curve at date $t$ ($M$
constant-maturity SVENY pillars) and $x_t\in\R$ the observed overnight short
rate. An encoder $E_\phi$ (a bidirectional LSTM) compresses a lookback window
into a latent initial state
\begin{equation}
z_0 = E_\phi\!\big(\mathbf c_{t-L:t},\,x_{t-L:t}\big)\in\R^d .
\end{equation}
The latent state evolves under a Neural SDE with learnable drift $\mu_\theta$
and diagonal diffusion $\sigma_\theta$,
\begin{equation}
\dd z_s = \mu_\theta(s,z_s)\,\dd s + \sigma_\theta(s,z_s)\,\dd W_s^{\QQ},
\qquad s\in[0,T_{\max}],
\label{eq:sde}
\end{equation}
integrated by Euler--Maruyama on a grid of step $\Delta=1/\texttt{spy}$. Two
drift families are supported,
\begin{equation}
\text{(simple)}\ \ \mu_\theta=f_\theta(s,z),
\qquad
\text{(OU)}\ \ \mu_\theta=\kappa_\theta(s,z)\big(\vartheta_\theta(s,z)-z\big),
\quad 0\le\kappa_\theta\le\kappa_{\max}.
\end{equation}
A decoder $D_\psi:\R^d\to\R$ produces the short rate with an additive anchor
pinning the front to the observed overnight rate,
\begin{equation}
r_s = D_\psi(z_s) + \big(x_t - D_\psi(z_0)\big)
\quad\Longrightarrow\quad r_0 = x_t .
\label{eq:anchor}
\end{equation}
Because all Monte-Carlo paths share $z_0$, the bracket in \eqref{eq:anchor} is
deterministic; consequently $r_s = x_t + \big(D_\psi(z_s)-D_\psi(z_0)\big)$,
i.e.\ the short-rate \emph{dynamics} are the decoder's response to the latent
displacement, anchored at $x_t$.

\section{Risk-neutral pricing}
\paragraph{Zero-coupon bonds and yields.}
\begin{equation}
P(0,\tau)=\E^{\QQ}\!\Big[e^{-\int_0^\tau r_s\,\dd s}\Big],
\qquad
y(\tau)=-\frac{\log P(0,\tau)}{\tau}.
\end{equation}
The expectation is a Monte-Carlo mean over $N$ paths, computed in log-space
for stability:
\begin{equation}
\log P(0,\tau)=\lse_{n}\big(-I^{(n)}_\tau\big)-\log N,
\qquad I^{(n)}_\tau=\Delta\!\!\sum_{s\le\tau} r^{(n)}_s .
\end{equation}
\paragraph{Convexity decomposition.}
A second-order cumulant expansion of $\log\E^{\QQ}[e^{-J}]$ with
$J=\int_0^\tau D_\psi(z_s)\,\dd s$ yields
\begin{equation}
y(\tau)\approx
\underbrace{(x_t-D_0)}_{\text{level}}
+\underbrace{\tfrac1\tau\!\int_0^\tau\!\E[D_\psi(z_s)]\dd s}_{\text{expectations}}
-\underbrace{\tfrac1{2\tau}\Var\!\Big(\!\int_0^\tau\! D_\psi(z_s)\dd s\Big)}_{\text{convexity}} .
\label{eq:convexity}
\end{equation}
The convexity term scales with the squared \emph{implied short-rate vol}
$\sigma_r\approx\lVert\nabla D_\psi\rVert\,\sigma_z$. A realistic
$\sigma_r\sim1\%$/yr produces $\mathcal O(20\,\text{bp})$ of $10$-year
convexity; an uncontrolled $\sigma_r\sim40\%$/yr produces
$\mathcal O(\text{hundreds of }\%)$, saturating the discount factor and
destroying the curve. This is the central reason volatility must be pinned
(\S\ref{sec:problems}, \S\ref{sec:terms}).
\paragraph{Treasury futures.}
For a contract with delivery $T_f$, deliverable basket $\mathcal B$ and
conversion factors $\{\mathrm{CF}_i\}$,
\begin{equation}
F=\E^{\QQ}\!\Big[\min_{i\in\mathcal B}\frac{B_i(z_{T_f})}{\mathrm{CF}_i}\Big],
\label{eq:futures}
\end{equation}
the expectation of the cheapest-to-deliver (CTD) ratio at delivery. The
forward deliverable-bond prices $B_i$ are produced by a learned
\emph{BondNet} $B_\chi(z,b_i)$ (bond features $b_i$), avoiding a nested
simulation inside \eqref{eq:futures}.

\section{Identifiability problems}\label{sec:problems}
Development proceeded by repeatedly finding that the calibration data does
not constrain quantities the model needs. The three structural ones:

\begin{enumerate}[leftmargin=*]
\item \textbf{Volatility is not identified.} Yields depend on $\sigma$ only
through the bp-level convexity term in \eqref{eq:convexity}; futures only
weakly, through CTD optionality. Left free, $\sigma$ drifts to whatever the
optimiser's side-pressures dictate. Empirically the trained model collapsed
to $\sigma_r\approx0.1$--$0.4\%$/yr and the Monte-Carlo fan degenerated to a
single line --- the model became effectively deterministic.

\item \textbf{The futures head was decoupled from the rates.} BondNet in
\eqref{eq:futures} is a free function of $z$; nothing tied $B_\chi$ to the
discount factors implied by the \emph{same} short-rate paths. Hence ``joint
calibration'' was two heads on a shared encoder, and a naive consistency
penalty whose gradient flowed into the SDE \emph{bent the yield curve to
meet BondNet's par-anchored prices}, producing a spurious
$-250$ to $-300$\,bp droop in the $5$--$10$y yields.

\item \textbf{The term premium is not identified under $\QQ$.} A pure
risk-neutral calibration fits prices but says nothing about the physical
dynamics; the market price of risk $\lambda$ is unobservable from the
cross-section of prices alone.
\end{enumerate}

\noindent Two further mechanism-level issues were diagnosed and fixed: (i) an
Ornstein--Uhlenbeck drift with unbounded mean-reversion rate $\kappa$ erases
the encoder signal $z_0$ over the horizon $1/\kappa$, collapsing the curve to
a constant (since $y(\tau)\propto g(\kappa\tau)$ with $g(x)=1-(1-e^{-x})/x\to1$
as $\kappa\to\infty$); (ii) scale and unit faults (percent vs.\ decimal
rates, a $365$-vs-$252$ day-count inconsistency between the calendar-year
maturity grid and date-derived horizons) that mis-timed the futures and
mis-scaled the encoder input.

\section{Resolving terms in the objective}\label{sec:terms}
The training objective is
\begin{equation}
\mathcal L=
\lambda_y\mathcal L_{\text{yield}}
+\lambda_f\mathcal L_{\text{fut}}
+\lambda_c\mathcal L_{\text{cons}}
+\lambda_\sigma\mathcal L_{\text{vol}}
+\lambda_{\PP}\mathcal L_{\PP/\QQ}.
\label{eq:objective}
\end{equation}

\paragraph{Yield and futures.}
$\mathcal L_{\text{yield}}=\frac1M\sum_\tau(y(\tau)-y^\star(\tau))^2$
(absolute, decimals). Futures use a \emph{relative} error so the weights are
commensurable, $\mathcal L_{\text{fut}}=\frac1{|\mathcal F|}\sum_k\big((F_k-F^\star_k)/F^\star_k\big)^2$;
absolute price MSE on $\$120$ contracts otherwise dwarfs the yield MSE by
$10^3$--$10^5$.

\paragraph{BondNet--SDE consistency (LSMC).}
Reconstruct each deliverable's cashflow schedule from its features and
discount it \emph{pathwise} on the model's own simulated rates,
$\hat B_i^{(n)}=\sum_j c_{ij}\,e^{-(I^{(n)}_{t_{ij}}-I^{(n)}_{T_f})}$. Then
\begin{equation}
\mathcal L_{\text{cons}}
=\frac1{100^2}\,\E_n\!\Big[\big(B_\chi(z^{(n)}_{T_f},b_i)-\sg[\hat B_i^{(n)}]\big)^2\Big],
\end{equation}
a Longstaff--Schwartz regression: $\E[\hat B_i\mid z_{T_f}]$ is the
model-consistent bond price, so $B_\chi$ learns it without nested
simulation. The stop-gradient $\sg[\cdot]$ is essential --- the gradient
reaches BondNet only; the non-detached variant is exactly what bent the
curve (\S\ref{sec:problems}.2).

\paragraph{Volatility anchor.}
Pin the model's $1$-year cross-path standard deviation to a historically
measured rate vol $\sigma^\star\approx1\%$,
\begin{equation}
\mathcal L_{\text{vol}}=\big(\std_n[r^{(n)}_{1\mathrm y}]-\sigma^\star\big)^2 .
\end{equation}
This is principled, not a fudge: by Girsanov the diffusion coefficient is
\emph{measure-invariant}, so a $\PP$-estimated vol is a valid prior for the
$\QQ$-dynamics.

\paragraph{$\PP/\QQ$ consistency and the term premium.}\label{par:pq}
The drift $\mu_\theta$ in \eqref{eq:sde} is the risk-neutral drift
$\mu^{\QQ}$. With market price of risk $\lambda\in\R^d$ and
$\dd W^{\QQ}=\dd W^{\PP}+\lambda\,\dd s$, Girsanov gives the physical drift
\begin{equation}
\mu^{\PP}(s,z)=\mu^{\QQ}(s,z)+\sigma_\theta(s,z)\,\lambda .
\label{eq:girsanov}
\end{equation}
A one-step physical forecast of the short rate is matched to the realised
future rate $h$ ahead,
\begin{equation}
\hat r^{\PP}_{t+h}=x_t+D_\psi\big(z_0+\mu^{\PP}(0,z_0)h\big)-D_\psi(z_0),
\qquad
\mathcal L_{\PP/\QQ}=\big(\hat r^{\PP}_{t+h}-x_{t+h}\big)^2 .
\end{equation}
This identifies $\lambda$ --- the term premium, since the instantaneous
forward equals the expected future rate plus $\sigma\lambda$. With
$\lambda_{\PP}=0$ the model is a pure $\QQ$-calibration and $\lambda\equiv0$;
the grid runs both ($\QQ$-only and joint $\PP/\QQ$) as a primary axis.

\paragraph{Numerical stabilisers.}
A short Euler unroll ($\Delta=1/32$), multiplicative coefficient scales
(setting $\sigma_r\approx1\%$ by construction rather than by a saturating
clamp), near-identity SDE initialisation, gradient clipping, a non-finite
gradient guard, and a BondNet output bias initialised near par ($\$100$)
together make the $\mathcal O(300)$-step backward pass stable.

\section{Identifiability summary}
\begin{center}
\begin{tabular}{@{}lll@{}}
\toprule
Quantity & Identified by the data? & Resolving term \\
\midrule
curve level / slope & yes (yields) & $\mathcal L_{\text{yield}}$ \\
short-rate volatility & \textbf{no} (vol-insensitive) & $\mathcal L_{\text{vol}}$ anchor \\
deliverable bond prices & only via CTD min & $\mathcal L_{\text{cons}}$ (LSMC) \\
term premium $\lambda$ & \textbf{no} under $\QQ$ alone & $\mathcal L_{\PP/\QQ}$ \\
\bottomrule
\end{tabular}
\end{center}
The volatility and term-premium terms are anchors/priors, \emph{not} fits;
this is a property of the data, not of the architecture.

\section{Open problems and future work}\label{sec:future}
\begin{enumerate}[leftmargin=*]
\item \textbf{The single-$z_0$ bottleneck.} All day-specific information
enters only through the initial condition; the drift, diffusion and decoder
are global. Conditioning $\mu_\theta$ (and, for the OU family, the long-run
mean $\vartheta_\theta$) on the encoder embedding would let the dynamics ---
not just the starting point --- be day-specific, and would in particular fix
the OU mean-reversion pathology at its root.
\item \textbf{One-factor decoder.} A linear decoder makes the short rate a
scalar projection $r=w^\top z+b$, i.e.\ a one-factor model regardless of $d$.
A multi-layer decoder (gridded) restores multi-factor curvature; whether the
data supports it is an empirical question the grid will answer.
\item \textbf{Honest model selection.} Selection should use a per-maturity
yield RMSE (in bp) on a held-out temporal fold, not the aggregate training
loss, whose scale is dominated by anchor/prior terms.
\item \textbf{A richer $\PP$ model.} The current $\PP/\QQ$ term is a local,
one-step drift match with a constant $\lambda$. A full treatment would make
$\lambda=\lambda(s,z)$ state-dependent and match the $\PP$-transition density
over the window (a likelihood, not a least-squares drift), yielding a
genuine joint estimation of the physical and risk-neutral dynamics and a
time-varying term premium.
\item \textbf{Arbitrage-consistent futures.} The LSMC term ties BondNet to
the model rates in expectation; a stronger guarantee would price
deliverables directly off the simulated curve, making the futures
exactly arbitrage-consistent with the yields by construction.
\end{enumerate}

\paragraph{Status.} The front of the curve fits to $\pm 25$\,bp, training is
stable, and futures price to $\sim 1\%$. The remaining failure mode under
investigation is a long-end drift droop; the detached LSMC term, the
volatility anchor, and the $\PP/\QQ$ axis introduced here target it directly,
with falsifiable signatures (a non-degenerate fan of width
$\sim\sigma^\star\sqrt\tau$, a flat per-maturity error, and a recovered
$\lambda$ consistent with a positive term premium).

\end{document}
